% Generated from validated Article Draft Result. Do not hand-edit; revise the structured draft instead.
\section{Runtimeは同じ月に何を吸収したか — Transformers v5・vLLM・SGLang}
\noindent\textbf{Transformers v5、vLLM、SGLangの1月更新を見ると、新model family、tool/API、diffusion、long-contextの変化が共通実装・serving層へ短い周期で流れ込んでいた。Agent化はruntime側にも継続的な再設計を要求する。}\autocite{src-152188dd686e,src-3bc80a3c7870,src-6a97353d35a9}

Hugging Faceは1月26日にTransformers v5をreleaseした。release notesは5年ぶりのmajor releaseとして、dynamic weight loading / WeightConverter、tokenizer changes、deprecated loading pathからのmigrationを含むbreaking changeと広いrefactoringを説明している。これは単なるmodel追加ではなく、共通実装層そのものの更新である。\autocite{src-152188dd686e}


\begin{claimboundary}
同日にv5.0.0rc3も記録され、GLM-4.7/GLM-ImageやMiniMax-M2のsupportが入ったが、stable v5を主要event、rc3をmodel-integrationの補助証拠として扱う。library supportはimplementation compatibilityを示すのであって、upstream modelの品質を評価したものではない。\autocite{src-152188dd686e,src-9f7ae4228266}
\end{claimboundary}

vLLM v0.14.0は1月20日、async schedulingをdefaultで有効化し、engine、tool calling、API、model supportを広く更新した。release notes自身がbreaking changeとし、unsupported configurationを明示している点は、serving最適化がdefault behaviorへ入る際の互換性コストも示している。\autocite{src-3bc80a3c7870}


9日後のv0.15.0ではKimi-K2.5 architectureがnew model supportへ加わり、async schedulingはpipeline parallelismへ拡張された。ここでKimi-K2.5はmodelそのものの性能記事ではなく、新しいarchitectureがserving runtimeへ取り込まれたdownstream uptakeの証拠として扱う。\autocite{src-d9f9de620203}


\begin{claimboundary}
vLLMのJanuary seriesは二つのreleaseを9日間隔で含み、一つのatomic featureではなくruntime evolutionとして読むべきだ。release notesのspeedup等はproject-reportedでconfiguration dependentであり、本稿では独立再現していない。\autocite{src-3bc80a3c7870,src-d9f9de620203}
\end{claimboundary}

SGLang v0.5.8も同月、GLM-4.7-Flashのday-zero support、diffusion serving、FlashAttention 4 decoding、long-context / disaggregated-servingの変更を取り込んだ。model、diffusion、long-contextという異質なworkloadがserving stackへ同時に流れ込んでいる。\autocite{src-6a97353d35a9}


\begin{claimboundary}
SGLangのreleaseは多数のPRと異なるhardware/model向けのperformance claimを束ねる。supportはimplementation availabilityの証拠であり、速度、TTFT、memoryの数値やproduction qualityを一般保証するものではない。\autocite{src-6a97353d35a9}
\end{claimboundary}

1月のruntime層は、新しいmodelを『動かせる』だけでなく、loading、tokenization、scheduling、tool/API、diffusion、long-contextを同時に更新していた。Agent execution stackが広がるほど、その下で共通runtimeが吸収すべき差分も増える。model releaseの速度とserving implementationの速度が同じ競争面へ近づき始めた。\autocite{src-152188dd686e,src-3bc80a3c7870,src-6a97353d35a9,src-d9f9de620203}
